\section{Methods}
\label{sec:methods}

\paragraph{Problem.}
Place $N$ turbines at positions $\mathbf{x}, \mathbf{y} \in
\mathbb{R}^N$ to maximize annual energy production
$\text{AEP}(\mathbf{x}, \mathbf{y}) = \sum_{d} w_d \sum_{i} P_i \cdot
8760$ subject to (i)~all positions inside a convex polygon
$\mathcal{P}$ and (ii)~pairwise distance $\geq s_{\min}$. Sums are
over wind direction--speed sectors $d$ with frequency weights
$w_d$, and over turbines $i$. The objective is non-convex (wake
interactions); the constraint set is convex.

\paragraph{Function interface.}
The LLM writes a single Python function
\texttt{optimize(sim, n\_target, boundary, min\_spacing, wd, ws,
weights)} returning $(\mathbf{x}, \mathbf{y})$, OR (in
\emph{schedule-only} mode) a single
\texttt{schedule\_fn(step, total\_steps, lr0, alpha0)} returning
$(\eta_t, \alpha_t, \beta_1, \beta_2)$ that controls a fixed
8000-step Adam skeleton. The skeleton handles initialization,
gradient computation, and constraint penalties; the LLM controls
\emph{only} the schedule. A harness loads the problem, instantiates
the wake model, calls the LLM's function, and scores the result.
The LLM cannot modify the wake model parameters.

\paragraph{Train/test protocol.}
The agent trains on a single cell: DEI polygon, $N\!=\!50$
turbines, observed DEI wind rose. This deliberate single-cell
design maximizes the strength of the generalization claim: any
positive transfer to other configurations represents
\emph{extrapolation}, not interpolation across a training
distribution.

We evaluate on a $2 \times 6 \times 4 = 48$-cell matrix
(Table~\ref{tab:matrix}) crossing two polygons, six turbine
counts, and four wind roses. The ROWP polygon is fully held
out---the agent never sees any ROWP configuration. The DEI
variants (different $N$, different wind roses) are also held out:
the agent sees only DEI-$N\!$50-DEI\_rose during training.

The LLM receives full AEP scores on the training cell; on
held-out, it sees only binary feasibility (PASS / FAIL). Held-out
AEP is logged silently for post-hoc analysis but never
communicated to the LLM.

\paragraph{Script selection.}
The \emph{deployed} schedule is the one with the highest training
AEP among scripts that pass all unit tests, including the
stressed-polygon feasibility test. This criterion uses no held-out
information and simulates a practitioner who lacks a validation
farm. We separately report oracle selection (best feasible
held-out AEP) as an upper bound. The same selection criterion is
applied to all models.

\paragraph{Baselines.}
We use two baselines, distinguished throughout by their multi-start
count.
\emph{(i) The 500-start \textsc{TopFarm} baseline} (the headline
number): 500 multi-start runs of
\texttt{topfarm\_sgd\_solve}~\cite{pedersen2019topfarm} with grid
initialization filtered to the polygon, 4000 gradient iterations with
Adam ($\beta_1=0.1$, $\beta_2=0.2$) preceded by 2000 constant-LR
iterations. The best feasible layout across 500 starts, gated at a
0.1\,m maximum-boundary-excursion tolerance (see below), is
5545.0\,GWh on DEI ($N{=}50$) and 4243.6\,GWh on ROWP ($N{=}74$).
We report best-\emph{feasible}, not best-overall: at strict-zero
tolerance the \textsc{TopFarm} solver is essentially infeasible
(3/500 feasible on DEI, 0/500 on ROWP), always landing millimetres to
centimetres outside the polygon, so a positive tolerance is required
for the baseline to exist at all.
\emph{(ii) The matched $K{=}50$ multi-start baseline} (used for the
per-cell schedule comparison): for a fair comparison, the seed
schedule and each discovered schedule are each run through the
\emph{same} fixed SGD skeleton from $K{=}50$ random restarts; each
restart is then filtered by the constraint tolerance $\tau$ (maximum
boundary excursion $\le\tau$ and minimum spacing $\ge d_{\min}-\tau$),
and the reported AEP is the maximum over the feasible restarts. This
isolates the effect of the learning-rate and penalty schedule, since
baseline and schedule share initialization, step count, and solver.

\paragraph{Constraint tolerance.}
Because the SGD schedule is designed to satisfy the boundary and
spacing constraints only to a tolerance near the terminal learning
rate---not exactly---we report the per-cell AEP comparison as a
function of a constraint tolerance $\tau$ rather than at a single
feasibility threshold. A layout is $\tau$-feasible if every turbine
lies within $\tau$ of the polygon and all pairwise spacings are at
least $d_{\min}-\tau$.

\paragraph{Unit tests.}
A three-tier sandbox test suite validates each submission:
signature check, 3-turbine quick run, and a stressed thin-rhombus
polygon with 25 turbines. The stressed-polygon test catches
optimizers that satisfy constraints on the spacious training
polygon but produce boundary violations on tight geometries; in
our experiments, every custom optimizer that later failed on
held-out ROWP also failed the stressed test, making it a 7-second
proxy for held-out feasibility.

\paragraph{Agent architecture.}
A tool-calling loop with six tools: \texttt{read\_file},
\texttt{write\_file}, \texttt{run\_tests}, \texttt{run\_optimizer}
(scores on training, returns AEP/time/feasibility),
\texttt{test\_generalization} (returns held-out PASS/FAIL only,
\emph{not} AEP), and \texttt{get\_status}. Generated code runs in a
restricted environment: no network, no API keys, filesystem writes
limited to the workspace, and a static blocklist on
\texttt{subprocess}/\texttt{exec}/\texttt{eval} within optimizer
code. The system prompt includes phase-2 nudges toward custom
optimizers after 30\% of the time budget elapses, and a diversity
nudge after 5 consecutive \texttt{topfarm\_sgd\_solve} wrappers.

\paragraph{Metrics.}
We report training AEP, feasible held-out AEP, held-out feasibility
rate, strategy classification, and time-to-first-improvement.
\emph{All ROWP scores in this paper are filtered to feasible
layouts only}---a critical filter, since several schedules
producing the highest \emph{nominal} ROWP scores violate
constraints (Sec.~\ref{sec:results}).
